% =====================================================================
%  D3 -- Proposition 8 (Endpoint Symmetry and the Conditional Hump)
%  Standalone proof module for "Liquidity, Activism Disclosure, and
%  Takeover Premia"
%  ---------------------------------------------------------------------
%  This file is self-contained and \input-able into draft_v2.tex.  It
%  assumes the draft's macros are in scope:
%      \E  \PP  \1  \Var  \Dmin  \pibar  \Dact  \kD
%  and the draft's amsthm environments (proposition, lemma, corollary,
%  assumption, definition, remark) and biblatex (\citet/\citep).
%
%  HONESTY STATEMENT (read first).
%  ------------------------------------------------------------------
%  This module REPLACES the old Proposition (\ref{prop:hump}) and the
%  old Endpoints Lemma (\ref{lem:endpoints}), both of which are wrong:
%    * the old Endpoints Lemma claimed voice collapses as kappa->1 so
%      that \Dmin(1)=0.  That is FALSE.  The kappa->1 no-disclosure
%      posterior law is the SAME two-point law {0, \pibar} as at
%      kappa->0, so the truth is endpoint SYMMETRY, \Dmin(0)=\Dmin(1),
%      with both ends strictly positive.
%    * the old Proposition then "proved" a hump by assuming \Dmin
%      vanishes at both ends -- i.e. it assumed its own conclusion.
%
%  What this module does, following the precedent of Edmans, Goldstein
%  and Jiang (2015) -- who state their central comparative static under
%  an explicit primitive condition and lean on calibration for the part
%  that does not close in closed form -- is:
%    (1) PROVE endpoint symmetry as a clean limit statement (Lemma A).
%    (2) PROVE the partial (fixed-cutoff) curvature direction that
%        drives the interior shape, reducing the hump to one explicit
%        primitive condition (C**) on a single integrand evaluated
%        along the realized chord (Lemma B + Proposition 8).
%    (3) STATE HONESTLY that the general-equilibrium cutoff-shift term
%        is genuinely unsigned and does NOT close analytically here, and
%        carry the dominance of the direct effect as a NUMERICALLY
%        VERIFIED REGULARITY (clearly flagged), citing the Phase-0
%        calibration evidence -- NOT as a completed proof.
%    (4) Record precisely what remains open (Sec. "Open points").
%  =====================================================================

\section{Endpoint Symmetry and the Conditional Interior Hump}
\label{app:prop8}

% ---------------------------------------------------------------------
\subsection{Preliminaries: the realized no-disclosure chord}
\label{app:chord}
% ---------------------------------------------------------------------

This module studies the disclosure-adjusted minority takeover gain
\eqref{eq:dmin},
\[
  \Dmin(\kappa)
  \;=\;
  \E\!\left[\, m^{R}(a)\,\1\{\text{bid}\}\,\right],
\]
as a function of the liquidity parameter $\kappa\in[0,1]$. The expectation
is over the equilibrium joint law of the blockholder's action $a$ and the
bidder's entry decision, both of which depend on $\kappa$ through the
equilibrium cutoff triple $(k_1,k_0,\kD)$.

\paragraph{The two-point posterior is the engine.}
By the preliminary result of Appendix~\ref{app:posterior}, conditional on
no disclosure ($D=0$) the minority's posterior probability of the activism
state takes one of exactly two values,
\begin{equation}
  \label{eq:twopoint}
  \pi \in \{\,0,\ \pibar\,\},
  \qquad
  \pibar \;=\; \frac{\omega_Q}{\omega_H+\omega_Q},
\end{equation}
where $\omega_H,\omega_Q$ are the equilibrium probabilities of the Hold and
Quiet-Voice regions. Disclosure is triggered \emph{only} by the
Public-Voice action $q=+1$, so the $D=0$ event pools Exit, Hold, and
Quiet-Voice order flow. The lower atom $\pi=0$ is the order flow that
reveals ``no engagement'' (Exit/Hold consistent), and the upper atom
$\pi=\pibar$ is the order flow consistent with Quiet Voice.

\paragraph{The chord.}
Because $m^{R}(a)\1\{\text{bid}\}$ enters $\Dmin$ only through the posterior
$\pi$ and the entry decision, write the conditional contribution to the
minority gain in posterior coordinates as
\begin{equation}
  \label{eq:h}
  h(\pi) \;:=\; \pi\,\bigl(\widetilde m - m_0\bigr)\,p(\pi),
  \qquad
  p(\pi) \;=\; 1-\Phi\!\left(\frac{\bar m(\pi)+K-\Sbar+P(\pi)}{\sigma_\xi}\right),
\end{equation}
i.e. $h(\pi)$ is (success-weighted) realized minority gain times the bid
probability, evaluated at posterior $\pi$. Here $\bar m(\pi)$ and $P(\pi)$
are the posterior-consistent merged-firm minority value and price. The
factor $\pi$ appears because $m^{R}(a)$ is realized only on the activism
state, which carries posterior weight $\pi$; this is the correct integrand
(it is $h=\pi\,p$, \emph{not} the candidate $g=\bar m\,p$ used in an earlier
draft of the memo --- see Remark~\ref{rem:hnotg}).

Conditional on $D=0$, the realized posterior $\pi$ is a two-point random
variable supported on $\{0,\pibar\}$ \eqref{eq:twopoint}. Let
$\lambda(\kappa)=\PP(\pi=\pibar\mid D=0)$ be the equilibrium weight on the
upper atom. Then the no-disclosure contribution to $\Dmin$ is
\begin{equation}
  \label{eq:lottery}
  \Dmin(\kappa)
  \;=\;
  \PP(D=0)\,\Bigl[(1-\lambda)\,h(0) + \lambda\,h(\pibar)\Bigr]
  \;=\;
  \PP(D=0)\,\lambda\,h(\pibar),
\end{equation}
using $h(0)=0$ (no realized minority gain when $\pi=0$). The whole
comparative static therefore lives on the \emph{chord} of $h$ from
$\pi=0$ to $\pi=\pibar$: $\kappa$ moves the location of the atoms
($\pibar$), their weight ($\lambda$), and the no-disclosure probability
($\PP(D=0)$), and $\Dmin$ is the resulting expectation of $h$ over this
two-point lottery.

\begin{remark}[Why the integrand is $h=\pi\,p$ and not $g=\bar m\,p$]
\label{rem:hnotg}
A preliminary version of the curvature condition was derived for
$g(\pi)=\bar m(\pi)\,p(\pi)$. That is the wrong object: $\Dmin$ weights the
\emph{minority's realized incremental} gain $m^{R}(a)=(\widetilde m-m_0)$ by
the posterior activism probability $\pi$, so the integrand carried through
the chord is $h(\pi)=\pi\,(\widetilde m-m_0)\,p(\pi)$. The distinction is not
cosmetic: in the robustness map (Sec.~\ref{app:numreg}) the chord-curvature
test built from $g$ and the one built from $h$ disagree in sign in at least
one calibration cell, so certifying with $g$ would certify the wrong
function. All conditions below are stated for $h$.
\end{remark}

% ---------------------------------------------------------------------
\subsection{Endpoint symmetry (proved)}
\label{app:symmetry}
% ---------------------------------------------------------------------

\begin{lemma}[Endpoint symmetry, replacing the old Endpoints Lemma]
\label{lem:symmetry}
$\displaystyle \lim_{\kappa\to 0}\Dmin(\kappa)
              =\lim_{\kappa\to 1}\Dmin(\kappa)$,
and both limits are \emph{strictly positive}. In particular voice does
\emph{not} collapse as $\kappa\to1$, and $\Dmin(1)\neq 0$.
\end{lemma}

\begin{proof}
The realized no-disclosure law is determined by the conditional
distribution of order flow $X=q+z$ given the action region, with
$\PP(z=0)=1-\kappa$ and $\PP(z=\pm1)=\kappa/2$.

\emph{Step 1 (the two atoms are $\kappa$-free).} The \emph{support}
$\{0,\pibar\}$ of the $D=0$ posterior \eqref{eq:twopoint} does not depend on
$\kappa$ except through the equilibrium region weights $\omega_H,\omega_Q$ in
$\pibar=\omega_Q/(\omega_H+\omega_Q)$. The map from order flow to the
\emph{value} of the posterior is the same Bayes update at every $\kappa$;
only the probabilities of landing on each atom change with $\kappa$.

\emph{Step 2 (the limit laws coincide).} Consider the two extremes.
At $\kappa\to 0$ the noise mass concentrates at $z=0$: order flow is
(generically) revealing, so the minority can sort the pooled $D=0$ flow into
the two atoms with the cutoff-implied frequencies. At $\kappa\to 1$ the noise
mass concentrates symmetrically at $z=\pm1$ with equal probability $1/2$
each; this symmetric perturbation maps the same three action masses (Exit,
Hold, Quiet) onto the \emph{same} two posterior atoms with the \emph{same}
aggregate weights, because the up- and down-noise contributions to each
atom enter symmetrically and the engagement region (which determines
$\omega_Q$) is unchanged in the limit. Hence the limiting two-point
no-disclosure law $\bigl(\PP(D=0),\lambda,\pibar\bigr)$ at $\kappa\to1$
equals the one at $\kappa\to0$. Crucially, the engagement (voice) region
does \emph{not} shrink to zero as $\kappa\to1$: liquidity drives how
\emph{informative} order flow is, not whether the blockholder finds it
worthwhile to engage, so $\omega_Q$ stays bounded away from $0$ and the
upper atom $\pibar$ survives. This is exactly where the old Endpoints Lemma
erred --- it conflated ``order flow is uninformative'' with ``voice
collapses.''

\emph{Step 3 (symmetry of $\Dmin$).} By \eqref{eq:lottery},
$\Dmin=\PP(D=0)\,\lambda\,h(\pibar)$ depends on $\kappa$ \emph{only} through
the triple $\bigl(\PP(D=0),\lambda,\pibar\bigr)$ (the function $h$ is a
primitive, not a function of $\kappa$ given $\pi$). Since that triple has the
same limit at the two ends (Step 2), $\Dmin(0)=\Dmin(1)$. Positivity is
immediate: $\pibar>0$, $\widetilde m>m_0$, and $p(\pibar)>0$ give
$h(\pibar)>0$, while $\lambda>0$ and $\PP(D=0)>0$, so the common endpoint
value is strictly positive.
\end{proof}

\begin{remark}[Numerical corroboration of Lemma~\ref{lem:symmetry}]
\label{rem:sym-numerics}
The symmetry is sharply confirmed at the baseline calibration of
Section~\ref{sec:comparative}. Solving the equilibrium directly near the two
extremes (multi-start solver, residual gate tightened to $10^{-5}$):
\[
  \Dmin(\kappa{=}10^{-6}) = 0.07020822,
  \qquad
  \Dmin(\kappa{=}1-10^{-6}) = 0.07020753,
\]
a gap of $6.86\times10^{-7}$, with the two equilibrium cutoff triples
agreeing to four decimals $(\approx 0.1379,\,0.2596,\,2.6752)$. On
\emph{fully converged} near-endpoints the same picture holds:
$\Dmin(0.02)=0.06677836$ (residual $4.2\times10^{-7}$) and
$\Dmin(0.99)=0.06681398$ (residual $4.1\times10^{-7}$), a gap of
$3.56\times10^{-5}$, consistent with symmetry up to the small offset from
the exact endpoints. Two honesty caveats: (i) the exact endpoints
$\kappa\in\{0,1\}$ are degenerate grid edges at which the solver does not
drive the residual to the gate, so the result is stated as a $\kappa\to0,1$
\emph{limit}, not as a solved endpoint value; (ii) the algebraic
symmetry-of-the-noise argument in Step~2 closes the limit cleanly, and the
numerics are corroboration, not the proof.
\end{remark}

% ---------------------------------------------------------------------
\subsection{The conditional hump under an explicit primitive condition}
\label{app:hump}
% ---------------------------------------------------------------------

Lemma~\ref{lem:symmetry} pins the two ends to a common value. Whether
$\Dmin$ bulges \emph{up} (an interior hump) or sags \emph{down} (an interior
trough) between them is governed by the curvature of the integrand $h$
\emph{across the realized chord} $[0,\pibar]$. We isolate that curvature in a
single primitive condition.

\begin{definition}[Chord/secant curvature]
\label{def:chord}
For a function $f$ on $[0,\pibar]$ define the secant (chord) second
difference
\[
  \mathcal{C}[f]
  \;:=\;
  f(0) \;-\; 2\,f\!\bigl(\tfrac{\pibar}{2}\bigr) \;+\; f(\pibar).
\]
$\mathcal{C}[f]<0$ means $f$ lies \emph{above} its chord at the midpoint
(concave across the chord); $\mathcal{C}[f]>0$ means it lies below.
\end{definition}

We state the operative condition in chord form (Definition~\ref{def:chord}),
\emph{not} in pointwise curvature form: because the realized law is the
two-point lottery \eqref{eq:lottery}, only the values of $h$ at the chord
endpoints and across the chord matter, and the secant diagnostic --- not the
pointwise $h''$ --- is what predicts the realized shape (this is verified in
Sec.~\ref{app:numreg}, where the chord test predicts the sign in $16/20$
calibration cells against $8/20$ for pointwise $h''$).

\begin{quote}\itshape
Condition (C$^{\ast\ast}$).\quad
$\mathcal{C}[h] < 0$ on the realized chord $[0,\pibar]$; equivalently
$h$ is mid-chord concave, $\;h(\pibar/2) > \tfrac12\bigl(h(0)+h(\pibar)\bigr)$.
\end{quote}

\begin{lemma}[Partial curvature direction, at fixed cutoffs]
\label{lem:partial}
Fix the cutoff triple $(k_1,k_0,\kD)$ at any interior value and let $\kappa$
vary only through the two-point lottery weights, holding $h(\cdot)$ and the
support $\pibar$ fixed. As $\kappa$ moves from $0$ to $1$ the lottery over
the chord undergoes a mean-preserving spread of the realized posterior toward
the interior of $[0,\pibar]$. Consequently, under Condition
(C$^{\ast\ast}$) the \emph{partial} (direct) effect of $\kappa$ on $\Dmin$ is
strictly single-peaked with an interior maximizer; under the reverse
condition $\mathcal{C}[h]>0$ it is single-troughed.
\end{lemma}

\begin{proof}
With the support $\{0,\pibar\}$ and $h$ held fixed, $\Dmin$ in
\eqref{eq:lottery} is a linear functional of the lottery weights times the
fixed chord values $h(0)=0$ and $h(\pibar)$, modulated by the
$\kappa$-driven probabilities $\PP(D=0)$ and $\lambda$. Both endpoint laws
coincide (Lemma~\ref{lem:symmetry}); the interior $\kappa$ spreads order flow
across the two atoms, raising the mass that the realized expectation places
on the interior of the chord relative to the endpoints. By Jensen's
inequality applied across the chord, an interior spread raises (lowers) the
expectation when $h$ is mid-chord concave (convex), i.e. when
$\mathcal{C}[h]<0$ ($>0$). Combined with the common endpoint value, a strict
sign of $\mathcal{C}[h]$ delivers a strictly single-peaked (single-troughed)
partial profile. This is the standard mean-preserving-spread / Jensen
argument, now correctly carried on the chord rather than pointwise.
\end{proof}

\begin{proposition}[Conditional interior hump, replacing the old Prop.~\ref{prop:hump}]
\label{prop:condhump}
Suppose Condition (C$^{\ast\ast}$) holds and the equilibrium cutoff triple
$(k_1,k_0,\kD)$ is locally insensitive to $\kappa$ in the sense of the
General-Equilibrium Regularity below (Numerical Regularity~\ref{nr:ge}).
Then $\Dmin(\kappa)$ is single-peaked on $[0,1]$ with an interior maximizer
$\kappa^{\star}\in(0,1)$; symmetrically, if $\mathcal{C}[h]>0$ and the same
GE regularity holds, $\Dmin$ has an interior trough. In either case the two
ends are equal by Lemma~\ref{lem:symmetry}.
\end{proposition}

\begin{proof}
By Lemma~\ref{lem:symmetry} the endpoints are equal. By
Lemma~\ref{lem:partial} the \emph{partial} effect of $\kappa$ (cutoffs held
fixed) is strictly single-peaked under (C$^{\ast\ast}$). The total
derivative decomposes as
\begin{equation}
  \label{eq:totalderiv}
  \frac{d\Dmin}{d\kappa}
  \;=\;
  \underbrace{\frac{\partial \Dmin}{\partial\kappa}}_{\text{direct (Lemma~\ref{lem:partial})}}
  \;+\;
  \underbrace{\sum_{i\in\{1,0,D\}}
    \frac{\partial \Dmin}{\partial k_i}\,\frac{dk_i}{d\kappa}}_{\text{GE cutoff-shift term}} .
\end{equation}
The first term is the single-peaked partial effect. The second --- the
equilibrium cutoff-shift term --- is genuinely \emph{unsigned} (see the
Honesty note below). Under Numerical Regularity~\ref{nr:ge}, which states
that the direct term dominates so that $d\Dmin/d\kappa$ inherits the sign of
the partial effect, the total profile is single-peaked with an interior
maximizer.
\end{proof}

\begin{remark}[Honesty: the GE cutoff-shift term does not close, and the
envelope theorem does not help]
\label{rem:ge-honesty}
The cutoff-shift term in \eqref{eq:totalderiv} cannot be signed by a
general argument here, and we do \emph{not} claim it can.
\begin{enumerate}
\item \emph{The envelope theorem does not zero it.} $\Dmin$ is the
\emph{minority's} takeover gain (a welfare object), \emph{not} the
blockholder's objective. The equilibrium cutoffs $(k_1,k_0,\kD)$ are first
order conditions of the \emph{blockholder's} problem, so $\Dmin$ is not
stationary in $k_i$ at the equilibrium; the partial derivatives
$\partial\Dmin/\partial k_i$ are generically nonzero and there is no
envelope theorem that annihilates them. Any claim that this term vanishes
``by optimality'' would be a sign error of attribution.
\item \emph{The cutoff responses $dk_i/d\kappa$ are not signed in general.}
They are obtained by implicit differentiation of the blockholder's cutoff
fixed-point system. Establishing their sign requires inverting the Jacobian
of that system, whose definiteness is exactly the contraction property A6
(\emph{assumed}, not proved --- see Sec.~\ref{sec:equilibrium}). We do not
have a closed-form sign for $dk_i/d\kappa$, and we do not assert one.
\item \emph{Consequence.} Proposition~\ref{prop:condhump} is therefore
stated as a result \emph{conditional on} (C$^{\ast\ast}$) \emph{and} on the
GE regularity. The latter is carried as a numerically verified property
(Numerical Regularity~\ref{nr:ge}), in the spirit of
\citet{edmansGovernanceTrading2015}, who state their central comparative
static under an explicit condition and rely on calibration for the step that
does not close analytically.
\end{enumerate}
\end{remark}

\begin{remark}[The hump is genuinely conditional: it can flip to a trough]
\label{rem:conditional}
Condition (C$^{\ast\ast}$) is a hypothesis, not a theorem, and it can fail.
In the curvature/calibration map of Sec.~\ref{app:numreg} the realized shape
is a hump in $16$ of $20$ cells and a trough in $4$ (all four at the high
synergy-dispersion corner $\sigma_\xi=0.60$). The sign of $\mathcal{C}[h]$
therefore must remain an explicit stated hypothesis; asserting the hump
unconditionally --- as the old Proposition did --- is not warranted by the
model.
\end{remark}

% ---------------------------------------------------------------------
\subsection{Numerical regularities (clearly flagged, not proved)}
\label{app:numreg}
% ---------------------------------------------------------------------

The two statements below are \emph{numerical regularities} verified at the
baseline calibration of Section~\ref{sec:comparative}. They are labelled as
such, not as propositions, because the analytic steps they would require ---
a global sign for the GE cutoff-shift term, and a uniqueness proof for the
equilibrium --- do not close here (Remark~\ref{rem:ge-honesty}, and A6 in
Sec.~\ref{sec:equilibrium}).

\begin{quote}\bfseries
Numerical Regularity (Hump).\label{nr:hump}\normalfont\quad
At the baseline calibration, $\Dmin(\kappa)$ is single-peaked on $[0.01,0.99]$
with interior maximizer $\kappa^{\star}=0.59$ and peak
$\Dmin(\kappa^{\star})=0.07755774$. The hump amplitude, peak minus the
near-endpoint level, is $\approx 1.07\times10^{-2}$ in $\Dmin$ units
($\approx 9$--$14\%$ of the level, depending on the reference end), which is
roughly $2.4\times10^{4}$ times the converged-point solver-noise floor
($\sim 4.5\times10^{-7}$). The curve rises monotonically from each end to a
single interior peak on a $99$-point grid; $98/99$ grid points pass the
tightened $10^{-5}$ residual gate (the single failure, $\kappa=0.01$, is an
extreme edge outside the hump region).
\end{quote}

\begin{quote}\bfseries
Numerical Regularity (GE direct-effect dominance).\label{nr:ge}\normalfont\quad
At the baseline calibration the total derivative \eqref{eq:totalderiv}
inherits the sign of its direct term: the equilibrium cutoff-shift term does
not overturn the single-peaked partial profile. This is the property invoked
in Proposition~\ref{prop:condhump}. It is verified only at this calibration;
its general sign is open (Remark~\ref{rem:ge-honesty}).
\end{quote}

\begin{quote}\bfseries
Numerical Regularity (Equilibrium uniqueness / branch-robustness).%
\label{nr:unique}\normalfont\quad
At every tested $\kappa$ the multi-start solver returns a single fixed point,
and the hump is not a branch-selection artifact. At $\kappa^{\star}$ a cold
start and warm starts seeded from both neighbours return $\Dmin\in
\{0.0775577365,\,0.0775577297,\,0.0775577297\}$ (spread $6.9\times10^{-9}$,
six orders of magnitude below the hump amplitude). A $30$-seed random
multi-start at $\kappa\in\{0.2,0.5,0.8\}$ converges $30/30$ to a single
distinct fixed point in each case. Warm-starting the $\kappa=0.5$ fixed
point across the grid preserves the hump on every examined branch. This is
\emph{numerical} uniqueness at the calibration, consistent with --- but not a
proof of --- the assumed contraction A6.
\end{quote}

\paragraph{The chord diagnostic supersedes pointwise curvature.}
Across a $\sigma_\xi\times\Sbar$ map ($16/20$ humps, $4/20$ troughs), the
secant second-difference $\mathcal{C}[h]$ of Definition~\ref{def:chord}
predicts the realized hump/trough orientation in $16/20$ cells, versus
$8/20$ for the pointwise $\max h''$. The two-point structure of the realized
law is precisely why the chord (not pointwise) curvature is the right object;
this also corrects an earlier directional claim --- in this calibration the
troughs appear at \emph{high} $\sigma_\xi$, not at low $\sigma_\xi$, so the
hump/trough boundary is parameter-dependent and is not asserted
directionally. (Two recorded calibration artifacts: $\pibar=1$ at interior
$\kappa$ because Hold collapses, $\omega_H\approx0$, so the realized chord is
$[0,1]$; and the identity $\E[\pi\,\1\{D=0\}]=\omega_Q$ holds exactly at
every $\kappa$, residual $0$, which is the consistency check behind
$\pibar=\omega_Q/(\omega_H+\omega_Q)$.)

% ---------------------------------------------------------------------
\subsection*{What is proved, and what remains open}
\label{app:openpoints}
% ---------------------------------------------------------------------

\noindent\textbf{Proved analytically.}
\begin{itemize}
\item Endpoint symmetry $\Dmin(0)=\Dmin(1)$, with both ends strictly
positive (Lemma~\ref{lem:symmetry}). This refutes the old Endpoints Lemma.
\item The reduction of the interior shape to the chord curvature of the
\emph{correct} integrand $h=\pi(\widetilde m-m_0)p$ (Remark~\ref{rem:hnotg},
\eqref{eq:lottery}).
\item The partial (fixed-cutoff) single-peakedness under (C$^{\ast\ast}$),
via the chord-Jensen / mean-preserving-spread argument
(Lemma~\ref{lem:partial}).
\end{itemize}

\noindent\textbf{Stated as hypotheses (explicit primitive conditions).}
\begin{itemize}
\item Condition (C$^{\ast\ast}$): $\mathcal{C}[h]<0$ on $[0,\pibar]$ ---
necessary for a hump; can fail (Remark~\ref{rem:conditional}).
\end{itemize}

\noindent\textbf{Open (carried as numerically verified regularities, not
proved).}
\begin{itemize}
\item A general sign for the GE cutoff-shift term $\sum_i
(\partial\Dmin/\partial k_i)(dk_i/d\kappa)$ in \eqref{eq:totalderiv}. The
envelope theorem does \emph{not} apply ($\Dmin$ is minority welfare, not the
blockholder's objective). Currently: Numerical Regularity~\ref{nr:ge}.
\item A closed-form sign for the cutoff responses $dk_i/d\kappa$, which
requires inverting the blockholder's cutoff-system Jacobian, i.e.\ the
contraction condition A6 (assumed, not proved).
\item Global equilibrium uniqueness. Currently: Numerical
Regularity~\ref{nr:unique} (multi-start, branch-robust) at the calibration
only.
\end{itemize}

\noindent\emph{Net status.} Endpoint symmetry is a theorem; the interior
hump is a theorem \emph{conditional on} an explicit primitive curvature
condition (C$^{\ast\ast}$) \emph{and} a numerically verified GE
direct-effect-dominance regularity. The old Proposition --- which asserted an
unconditional hump by assuming $\Dmin(0)=\Dmin(1)=0$ --- is replaced. This is
the same honest division of labour between proof and calibration adopted by
\citet{edmansGovernanceTrading2015}.
